\documentclass[10pt,a4paper]{amsart}
\usepackage[margin=19mm]{geometry}
\usepackage{amsmath,amssymb,mathtools}
\usepackage[scr=rsfs,cal=euler]{mathalfa}
\usepackage{xfrac}
\usepackage[unicode,hidelinks,bookmarks=false]{hyperref}
\newcommand{\ie}{i.e.\ }
\newcommand{\cf}{cf.\ }
\newcommand{\resp}{respectively\ }
\newcommand{\Subsection}{\subsection}
\providecommand{\nobreakdash}{\nobreak\hbox{-}}
\setlength{\emergencystretch}{2em}
\multlinegap0pt
\usepackage{bm}
\usepackage{bbm}
\usepackage{dsfont}
\usepackage{xspace}
\usepackage{cancel}
\usepackage{mathtools}
\usepackage{amsthm}
\makeatletter
\@ifundefined{theorem}{\newtheorem{theorem}{Theorem}}{}
\@ifundefined{lemma}{\newtheorem{lemma}[theorem]{Lemma}}{}
\makeatother
\newcommand {\Sa}{\mathbb{S}^{\sharp_{A}}}
\newcommand {\Soa}{S_{1}^{\sharp_{A}}}
\newcommand {\Sta}{S_{2}^{\sharp_{A}}}
\newcommand {\Ta}{\mathbb{T}^{\sharp_{A}}}
\newcommand {\Toa}{T_{1}^{\sharp_{A}}}
\newcommand {\Tta}{T_{2}^{\sharp_{A}}}
\newcommand {\T}{\mathbb{T}}
\newcommand {\Wa}{\mathbb{W}^{\sharp_{A}}}
\newcommand {\W}{\mathbb{W}}
\newcommand {\Za}{\mathbb{Z}^{\sharp_{A}}}
\newcommand {\Z}{\mathbb{Z}}
\newcommand {\h}{\mathcal{H}}
\newcommand{\abs}[1]{\left\vert#1\right\vert}
\newcommand{\bra}[1]{\left(#1\right)}
\newcommand{\hh}{\mathcal{H}\oplus\mathcal{H}}
\newcommand{\normA}[1]{{\left\Vert#1\right\Vert}_{A}}
\newcommand{\norm}[1]{\left\Vert#1\right\Vert}
\newcommand{\seqA}[1]{{\left<#1\right>}_{A}}
\renewcommand {\S}{\mathbb{S}}
\renewcommand {\b}{\mathcal{B}}
\title[Tighter Inequalities for $A$-Numerical Radii of Operator Mat]{Tighter Inequalities for $A$-Numerical Radii of Operator Matrices and Their Applications}
\author{M.H.M. Rashid}
\date{Source: arXiv:2507.05105v1 (CC BY 4.0)}
\begin{document}
\maketitle

\noindent\emph{Source and licence.} Tighter Inequalities for $A$-Numerical Radii of Operator Matrices and Their Applications. Version arXiv:2507.05105v1 (CC BY 4.0), distributed under the Creative Commons Attribution 4.0 International licence, as recorded in the arXiv licence metadata for this exact version and shown on the abstract page https://arxiv.org/abs/2507.05105v1. Full rights notice: licenses/arxiv-2507.05105-CC-BY-4.0.txt.\par\smallskip

\noindent\textit{Editorial scope.} Counted unit: the Theorem whose source label is \texttt{Thm:prod1} in the article (source0.tex lines 966--973, source label \texttt{Thm:prod1}), delivered together with the article's own complete original proof of it (source0.tex lines 974--1025, one \texttt{proof} environment of 52 lines). No mathematics is written, repaired or re-derived: statement and proof are the article's own text line for line. Required context retained uncounted: lem:holder\_mccarthy (lines 302--309); lem:jensen (lines 311--317); Buzano (lines 596--600). Every definition, display and label of those lines is retained unchanged. Omitted from this package and not redistributed: 1--301, 310--310, 318--595, 601--965, 1026--1682. Nothing inside a retained excerpt is omitted or altered: every retained body is delivered exactly as the article's lines are written, so no omission receipt of any kind accompanies this package. Citations: the retained excerpts carry no citation command at all, so no bibliography entry of the article is transcribed and no reference is redistributed. Reference adaptation: none was needed and none was made; every reference inside the delivered unit resolves inside it, so no unresolved reference or citation is produced. Printed slips kept literal: no printed word, symbol, hypothesis or number of the counted statement or of its proof was altered. Layout and class: the article's own class is \texttt{article} with packages including arxiv, inputenc, fontenc, hyperref, url, booktabs, amsfonts, nicefrac, microtype, lipsum, graphicx, natbib, doi, amsmath; the wrapper is the lane's pinned \texttt{amsart} editorial wrapper at 10pt on A4, which restates the theorem-like environments the retained text needs and adds the article's own macro definitions that the retained text needs (\texttt{\textbackslash S}, \texttt{\textbackslash Sa}, \texttt{\textbackslash Soa}, \texttt{\textbackslash Sta}, \texttt{\textbackslash T}, \texttt{\textbackslash Ta}, \texttt{\textbackslash Toa}, \texttt{\textbackslash Tta}, \texttt{\textbackslash W}, \texttt{\textbackslash Wa}, \texttt{\textbackslash Z}, \texttt{\textbackslash Za}, \texttt{\textbackslash abs}, \texttt{\textbackslash b}, \texttt{\textbackslash bra}, \texttt{\textbackslash h}, \texttt{\textbackslash hh}, \texttt{\textbackslash norm}, \texttt{\textbackslash normA}, \texttt{\textbackslash seqA}), with extra wrapper packages (bm, bbm, dsfont, xspace, cancel, mathtools). The delivered numbering is the unit's own. Assets: no retained span contains a figure environment, a graphics-inclusion command, a PDF-page inclusion, a table or a TikZ picture; no image file is copied into this package, the package declares no asset dependency and no image file is redistributed with it. Licence: version arXiv:2507.05105v1 of this article is distributed under the Creative Commons Attribution 4.0 International licence, as recorded in the arXiv licence metadata for this exact version and shown on the abstract page https://arxiv.org/abs/2507.05105v1. A full rights notice is delivered at licenses/arxiv-2507.05105-CC-BY-4.0.txt. Characters: every retained excerpt is pure ASCII, so the delivered engine, XeLaTeX with its default Latin Modern text font, reports no missing character.
\begin{lemma}[H\"older-McCarthy Inequality]
\label{lem:holder_mccarthy}
For $T \in \mathcal{B}_A(\mathcal{H})$ and $x \in \mathcal{H}$ with $\norm{x}_A = 1$:
\begin{enumerate}
    \item If $T$ is $A$-positive and $r \geq 1$, then $\langle Tx, x \rangle_A^r \leq \langle T^rx, x \rangle_A$.
    \item If $T$ is $A$-positive and $0 \leq r \leq 1$, then $\langle T^rx, x \rangle_A \leq \langle Tx, x \rangle_A^r$.
\end{enumerate}
\end{lemma}

\begin{lemma}[Jensen-Type Inequality]
\label{lem:jensen}
For $a, b > 0$ and $\alpha \in [0, 1]$, the following inequalities hold for any $r \geq 1$:
\[
a^\alpha b^{1-\alpha} \leq \alpha a + (1-\alpha)b \leq \left(\alpha a^r + (1-\alpha)b^r\right)^{1/r}.
\]
\end{lemma}

\begin{lemma}\label{Buzano} If $a,b,e\in\h$ with $\normA{e}=1$ and $\beta\geq 0$, then
\begin{equation}\label{Imediae1}
  \abs{\seqA{a,e}\seqA{e,b}}^2\leq \frac{1}{4}\bra{\frac{2\beta+1}{\beta+1}\normA{a}^2\normA{b}^2+\frac{2\beta+3}{\beta+1}\normA{a}\normA{b}\abs{\seqA{a,b}}}.
\end{equation}
\end{lemma}

 \begin{theorem}\label{Thm:prod1} Let $\T=\begin{bmatrix} O& T_1\\ T_2& O \\\end{bmatrix}, \S=\begin{bmatrix} O& S_1\\ S_2& O \\\end{bmatrix}\in\b_A(\hh)$ and $\beta\geq 0$. Then
 \begin{equation}\label{Eq:prod1}
 \omega_A^{4r}\bra{\Sa \T}\leq \frac{1+2\beta}{16(\beta+1)}\max\{\phi_r^2,\psi_r^2\}+\frac{3+2\beta}{8(\beta+1)}\max\{\phi_r,\psi_r\}
 \max\{\rho_r,\theta_r\}
 \end{equation}
 for all $r\geq 1$, where $\phi_r=\normA{\bra{\Tta T_2}^{2r}+\bra{\Sta S_2}^{2r}},\psi_r=\normA{\bra{\Toa T_2}^{2r}+\bra{\Soa S_2}^{2r}}$,
 $\rho_r=\omega_A^{r}\bra{\Sta S_2\Tta T_2}$ and $\theta_r=\omega_A^{r}\bra{\Soa S_1\Toa T_1}$.
 \end{theorem}

 \begin{proof} We assume that $\gamma_1=\frac{1+2\beta}{\beta+1}$,$\gamma_2=\frac{3+2\beta}{\beta+1}$,
 $(\Ta\T)^{2r}+(\Sa\S)^{2r}=$ \\$\begin{bmatrix} \bra{\Tta T_2}^{2r}+\bra{\Sta S_2}^{2r}& O\\ O&
 \bra{\Toa T_2}^{2r}+\bra{\Soa S_2}^{2r} \\\end{bmatrix}$
  and $\Sa \S\Ta\T=\begin{bmatrix}\Sta S_2\Tta T_2& O\\ O& \Soa S_1\Toa T_1 \\\end{bmatrix}$.

 Now, let $z\in\hh$ with $\normA{z}=1$, we have
 \begin{eqnarray*}
 % \nonumber % Remove numbering (before each equation)
   \abs{\seqA{\Za\W z,z}}^{4} &=&\abs{\seqA{\W z,z}\seqA{z,\Z z}}^2 \\
    &\leq& \frac{\gamma_1}{4}\normA{\W z}^2\normA{\Za z}^2+\frac{\gamma_2}{4}\normA{\W z}\normA{\Za z}\abs{\seqA{\W z,\Za z}}\\
    &&\bra{\mbox{by Lemma \ref{Buzano}}}\\
    &\leq& \bra{\frac{\gamma_1}{4}\normA{\W z}^{2r}\normA{\Za z}^{2r}
    +\frac{\gamma_2}{4}\normA{\W z}^{r}\normA{\Za z}^{r}\abs{\seqA{\W z,\Za z}}^{r}}^{\frac{1}{r}}\\
    &&\bra{\mbox{by Lemma \ref{lem:jensen}}}
 \end{eqnarray*}
 Consequently,
 \begin{eqnarray*}
 % \nonumber % Remove numbering (before each equation)
   \abs{\seqA{\Za\W z,z}}^{4r} &\leq&\frac{\gamma_1}{4}\bra{\sqrt{\seqA{\Wa\W z,z}^{r}\seqA{\Z\Za z,z}^{r}}}^{2}\\
   &+&
    \frac{\gamma_2}{4}\sqrt{\seqA{\Wa\W z,z}^{r}\seqA{\Z\Za z,z}^{r}}\abs{\seqA{\Z\W z, z}}^{r}\\
    &\leq&\frac{\gamma_1}{16}\bra{\seqA{\Wa\W z,z}^{r}+\seqA{\Z\Za z,z}^{r}}^2\\
    &+&\frac{\gamma_2}{8}\bra{\seqA{\Wa\W z,z}^{r}+\seqA{\Z\Za z,z}^{r}}\abs{\seqA{\Z\W z, z}}^{r}\\
    &&\bra{\mbox{by the arithmetic-geometric mean inequality}}\\
    &\leq& \frac{\gamma_1}{16}\bra{\seqA{\bra{\Wa\W}^{r} z,z}+\seqA{\bra{\Z\Za}^{r} z,z}}^2\\
    &+&\frac{\gamma_2}{8}\bra{\seqA{\bra{\Wa\W}^{r} z,z}+\seqA{\bra{\Z\Za}^{r} z,z}}\abs{\seqA{\Z\W z, z}}^{r}\\
    &&\bra{\mbox{by Lemma \ref{lem:holder_mccarthy}}}.
 \end{eqnarray*}
 By replacing $\W=\Ta \T$ and $\Z=\Sa \S$ in the above inequality, and using the fact that $\bra{\Ta \T}^{\sharp_A}=\Ta\T$ and
 $\bra{\Sa\S}^{\sharp_A}=\Sa\S$. Thus, it can be deduced that
  \begin{eqnarray*}
  % \nonumber % Remove numbering (before each equation)
    \abs{\seqA{\W z,z}\seqA{z,\Z z}}^{2r}&\leq&  \frac{\gamma_1}{16}\bra{\seqA{\bra{\bra{\Ta\T}^{2r}+\bra{\Sa\S}^{2r}} z,z}}^2\\
     \\
    &+&\frac{\gamma_2}{8}\seqA{\seqA{\bra{\bra{\Ta\T}^{2r}+\bra{\Sa\S}^{2r}} z,z}}\abs{\seqA{\Sa\S\Ta\T z, z}}^{r}
    \end{eqnarray*}
    \begin{eqnarray*}
    &\leq&\frac{\gamma_1}{16}\normA{\bra{\Ta\T}^{2r}+\bra{\Sa\S}^{2r}}^2\\
    &+& \frac{\gamma_2}{8}\normA{\bra{\Ta\T}^{2r}+\bra{\Sa\S}^{2r}}
    \omega_A^{r}\bra{\Sa\S\Ta\T}.
  \end{eqnarray*}
  In addition, by utilizing the Cauchy-Schwarz inequality, we get
  \begin{eqnarray*}
  % \nonumber % Remove numbering (before each equation)
    \abs{\seqA{z,\Sa\T z}}^{4r} &=&\abs{\seqA{\T z,\S z}}^{4r} \leq  \normA{\T z}^{4r}\normA{\S z}^{4r}\\
     &=& \seqA{\T z,\T z}^{2r}\seqA{\S z,\S z}^{2r}=\abs{\seqA{\Ta\T z,z}\seqA{\Sa\S z,z}}^{2r}.
  \end{eqnarray*}
  Combining above two inequalities, we can obtain
  $$\abs{\seqA{z,\Sa\T z}}^{4r} \leq \frac{\gamma_1}{16}\max\{\phi_r^2,\psi_r^2\}+\frac{\gamma_2}{8}\max\{\phi_r,\psi_r\}
  \max\{\rho_r,\theta_r\}.$$
  Taking the supremum over all vectors $z\in\hh$ with $\normA{z}=1$, we obtain the desired inequality.
  \end{proof}

\end{document}
